\section{FFT and convolution}
\label{sec:fft}

The FFT family is new in 1.1.0. Its mechanisms are implemented and tested; its
performance rows are pending the sweep, and Table \ref{tab:fft} says so in every
cell that would carry a number.

\subsection{Stockham, and why there is no bit reversal pass}

The transform is Stockham autosort~\cite{vanloan1992}. Each stage writes to a
second buffer with the digit permutation folded into the output address, so every
stage reads and writes contiguous runs. At stride $p$ a thread owns one butterfly:
it reads elements $j$ and $j + n/2$ of the source, twists the second by
$\exp(-2\pi i k / 2p)$ with $k = j \bmod p$, and writes the sum and the difference
to their permuted destinations. The radix 4 and radix 8 stages are the same
statement with $R$ in place of 2, and there is exactly one implementation of each
in the tree, called from two translation units, because a radix 8 DFT written
twice is a radix 8 DFT that will eventually disagree with itself.

The in place alternative needs a separate bit reversal pass whose access pattern is
a scatter across the whole array, one sector per element at large $n$. On a family
that is already memory bound that pass costs a full memory round trip, which is a
whole extra pass out of the five the four step form needs. The price of avoiding it
is the second buffer, and the shared memory budget pays it. Cooley-Tukey with a bit
reversal pass is not implemented here and is not planned.

\subsection{The four step form}

For $n$ above the shared resident bound the plan factors $n = n_1 n_2$ with both
factors at or below $2^{12}$ and runs a transpose, batched size $n_2$ transforms
with the cross twiddle folded into the store epilogue, a transpose, batched size
$n_1$ transforms, and a final transpose that puts the output where the caller wants
it. Five passes, each one read and one write, so $5 \times 16n$ bytes; against
$\log_2 n$ passes for the naive multi kernel rung that is 5 against 24 at
$n = 2^{24}$. The cross twiddle costs no pass of its own because it rides on an
epilogue, and neither does a fused convolution's pointwise multiply.

The factorization splits $\log_2 n$ as evenly as it can, larger half first, which
keeps both factors inside the shared resident bound for every $n$ up to $2^{24}$
and gives both batched passes the most blocks to fill 48 SMs with. An override
exists and is validated rather than trusted.

\subsection{Convolution, and the crossover}

The headline figure of this family is a crossover, not a percentage. The FFT
path's cost barely moves with the filter length $M$; the direct path is $2NM$ FLOP
and does not stop growing. Four variants are implemented: an FFT path with a
standalone pointwise multiply kernel, an FFT path with the filter spectrum multiply
folded into a store epilogue (the same epilogue fusion argument the GEMM bias study
uses, and the same one FlashAttention makes at a larger scale~\cite{dao2022}), a
tiled time domain kernel, and the same tiling with the taps in constant memory.

Two implementation notes are worth keeping. The obvious direct tiling, a signal
tile plus a halo of $M-1$ samples, does not survive the top of the sweep:
$M = 16384$ is a 64 KB halo before the tile is counted. Chunking the filter instead
holds the shared footprint at about 3 KB whatever $M$ is. And the constant memory
bank is one per module rather than one per plan, so a plan binds its taps once and
records that it owns them; a call from a plan that does not own the bank rebinds
first, and a direct call from an owner that never bound throws rather than
convolving with somebody else's filter.

The crossover chart itself is generated from the committed sweep rows, with the
crossing read out of the data by log interpolation and printed into the caption. It
is not drawn by hand, and when the summary holds no convolution rows the generator
removes the figure and writes the paragraph below instead of rebuilding this
chapter around a stale picture.

The convolution crossover chart is pending: the committed summary carries no convolution rows, so there is nothing to plot and nothing is drawn. It appears here, with the crossing filter length read out of the data, once the sweep has run the convolution grid.


The plan's own selection rule divides the two traffic models by the two roofs of
this part, the direct model by the FP32 roof and the FFT model by the measured DRAM
roof. That is arithmetic and not a measurement. The crossing it predicts is exactly
what the sweep exists to check, and the rule is provisional until the sweep runs.

\begin{table}[htbp]
\centering
\begin{tabular}{lllrr}
\toprule
rung & mechanism & traffic model & GFLOP/s at the largest swept size & percent of cuFFT \\
\midrule
radix 2, one kernel per stage & a launch per butterfly stage & log2(n) passes & 180.3 & 12.8 \\
shared resident & one block per transform, all stages in shared & 1 pass & 584.7 & 34.3 \\
radix 4 shared resident & half the stages of radix 2 & 1 pass & 359.8 & 25.5 \\
radix 8 shared resident & a third of the stages of radix 2 & 1 pass & 541.0 & 38.4 \\
four step & transpose, batched, transpose, batched, transpose & 5 passes & 473.4 & 33.5 \\
convolution, FFT separate & standalone pointwise multiply kernel & 24 bytes per point & 501.4 & 47.0 \\
convolution, FFT fused & multiply folded into a store epilogue & 8 bytes per point & 532.0 & 49.8 \\
convolution, direct shared & tiled time domain, filter chunked through shared & 2 N M FLOP & 2436.4 & 3.3 \\
convolution, direct constant & the same tiling, taps in constant memory & 2 N M FLOP & 3758.5 & 20.3 \\
\bottomrule
\end{tabular}

\caption{The FFT and convolution rungs. Traffic models are derived and stated in
\texttt{docs/fft.md}; measured throughput and percent of cuFFT are pending.}
\label{tab:fft}
\end{table}
